% SEGMENT OLP-0694-S01
% Part: proof-theory
% Chapter: propositions-as-types
% Section: sequent-natural-deduction

\documentclass[../../../include/open-logic-section]{subfiles}

\begin{document}

\olfileid{int}{pty}{snd}

\olsection{தொடரணிகளுடன் இயல்பான வருவித்தல்}

$\Gamma \Sequent !A$ என்று எழுதுவோம்: அதாவது,
$\Gamma$ என்பது !!{undischarged} எடுகோள்களின்
கணமாகவும் $!A$ முடிவாகவும் உள்ள இயல்பான வருவித்தல்
!!{derivation} ஒன்று உள்ளது. $\Gamma$ காலிக் கணம்
எனில், $\Sequent !A$ என்று எழுதுவோம்.

$\Gamma \cup \{!A_1, \dots, !A_n\}$ என்பதற்கு
$\Gamma, !A_1, \dots, !A_n$ என்றும்,
$\Gamma \cup \Delta$ என்பதற்கு $\Gamma, \Delta$
என்றும் எழுதுவோம்.

% SEGMENT OLP-0694-S02
$\Gamma \Sequent !A \land !B$ எனில்,
$\Gamma$ என்பதை !!{undischarged} எடுகோள்களாகவும்
$!A \land !B$ ஐ இறுதி-!!{formula} ஆகவும் கொண்ட
!!a{derivation} நமக்கு உள்ளது. அதன் அடியில்
$\Elim{\land}$ விதியைப் பயன்படுத்தி, அதே
!!{undischarged} எடுகோள்களிலிருந்து $!A$ ஐ முடிவாகக்
கொண்ட !!a{derivation} ஐப் பெறலாம். வேறு சொற்களில்,
$\Gamma \Sequent !A \land !B$ எனில்
$\Gamma \Sequent !A$:
\begin{prooftree}
  \Axiom$\Gamma \fCenter !A \land !B$
  \RightLabel{$\Elim{\land}$}
  \UnaryInf$\Gamma \fCenter !A$
  \DisplayProof\qquad\bottomAlignProof
  \Axiom$\Gamma \fCenter !A \land !B$
  \RightLabel{$\Elim{\land}$}
  \UnaryInf$\Gamma \fCenter !B$
\end{prooftree}
$\Elim{\land}$ என்ற அடையாளம், இயல்பான வருவித்தலில்
அதே பெயருடைய விதியுடனான தொடர்பைக் காட்டுகிறது.

% SEGMENT OLP-0694-S03
அதேபோல், $\Gamma, !A \Sequent !B$ எனக் கொள்வோம்.
அதாவது, $\Gamma, !A$ என்பவற்றை !!{undischarged}
எடுகோள்களாகவும் $!B$ ஐ இறுதி-!!{formula} ஆகவும்
கொண்ட !!a{derivation} உள்ளது. இதற்கு
$\Intro{\lif}$ விதியைப் பயன்படுத்தினால்,
$\Gamma$ என்பதை !!{undischarged} எடுகோள்களாகவும்
$!A \lif !B$ ஐ இறுதி-!!{formula} ஆகவும் கொண்ட
!!a{derivation} கிடைக்கும்; அதாவது,
$\Gamma \Sequent !A \lif !B$.
எடுகோள்களின் !!{discharge} இங்கு மேலும் தெளிவாகக்
காட்டப்படுவதைக் கவனிக்கவும்.
\begin{prooftree}
  \Axiom $\Gamma, !A \fCenter !B$
  \RightLabel{$\Intro{\lif}$}
  \UnaryInf $\Gamma \fCenter !A \lif !B$
\end{prooftree}

இதே முறையில் மற்ற விதிகளிலிருந்தும் முடிவுகளைக்
காணலாம்; அவை முழுவதும் பின்வருமாறு:
\begin{gather*}
  \Axiom $\Gamma \fCenter !A$
  \Axiom $\Delta \fCenter !B$
  \RightLabel{$\Intro{\land}$}
  \BinaryInf $\Gamma,\Delta \fCenter !A \land !B$
  \DisplayProof\\
  \Axiom $\Gamma \fCenter !A \land !B$
  \RightLabel{$\Elim{\land}_1$}
  \UnaryInf $\Gamma \fCenter !A$
  \DisplayProof
  \qquad
  \Axiom $\Gamma \fCenter !A \land !B$
  \RightLabel{$\Elim{\land}_2$}
  \UnaryInf $\Gamma \fCenter !B$
  \DisplayProof
  \\
  \Axiom $\Gamma \fCenter !A$
  \RightLabel{$\Intro{\lor}_1$}
  \UnaryInf $\Gamma \fCenter !A \lor !B$
  \DisplayProof\qquad
  \Axiom $\Gamma \fCenter !B$
  \RightLabel{$\Intro{\lor}_2$}
  \UnaryInf $\Gamma \fCenter !A \lor !B$
  \DisplayProof\\
  \Axiom $\Gamma \fCenter !A \lor !B$
  \Axiom $\Delta, !A \fCenter !C$
  \Axiom $\Delta', !B \fCenter !C$
  \RightLabel{$\Elim{\lor}$}
  \TrinaryInf $\Gamma, \Delta, \Delta' \fCenter !C$
  \DisplayProof
  \\
  \Axiom $\Gamma, !A \fCenter !B$
  \RightLabel{$\Intro{\lif}$}
  \UnaryInf $\Gamma \fCenter !A \lif !B$
  \DisplayProof
  \qquad
  \Axiom $\Delta \fCenter !A \lif !B$
  \Axiom $\Gamma \fCenter !A$
  \RightLabel{$\Elim{\lif}$}
  \BinaryInf $\Gamma, \Delta \fCenter !B$
  \DisplayProof
  \\
  \Axiom $\Gamma \fCenter \lfalse$
  \RightLabel{$\FalseInt$}
  \UnaryInf $\Gamma \fCenter !A$
  \DisplayProof
\end{gather*}

% SEGMENT OLP-0694-S04
ஓர் எடுகோள் தனியே $!A$ இலிருந்து $!A$ இன்
!!a{derivation} ஆகும்; எனவே $!A \Sequent !A$
எப்போதும் கிடைக்கும்.

\begin{prooftree}
  \AxiomC{$ $}
  \UnaryInfC{$!A \fCenter !A$}
\end{prooftree}

இவ்விதிகளை ஒருங்கே எடுத்தால், எந்த இயல்பான வருவித்தல்
!!{derivation}s உள்ளன என்பதை விவரிக்கும் கலனமாகக்
கொள்ளலாம். இயல்பான வருவித்தலின் குறியீட்டு மாற்றுருவாகவும்
கொள்ளலாம்: ஒவ்வொரு படியும் !!{derive} செய்யப்பட்ட
!!{formula} ஐ மட்டுமன்றி, அது எந்த !!{undischarged}
எடுகோள்களிலிருந்து !!{derive} செய்யப்பட்டது என்பதையும்
பதிவுசெய்கிறது.

எடுத்துக்காட்டாக, பின்வரும் நிறுவலைக் காண்க:
\begin{prooftree}
  \Axiom$!A \fCenter !A$
  \UnaryInf $!A \fCenter !A \lor (!A \lif \lfalse)$
  \Axiom $!B \fCenter !B$
  \BinaryInf$!A, !B \fCenter \lfalse$
  \UnaryInf$!B \fCenter !A \lif \lfalse$
  \UnaryInf$!B \fCenter !A \lor (!A \lif \lfalse)$
  \Axiom$!B \fCenter !B$
  \BinaryInf$!B \fCenter \lfalse$
  \UnaryInf$\fCenter !B \lif \lfalse$
\end{prooftree}
இங்கு $!B$ என்பது
$(!A \lor (!A \lif \lfalse))\lif \lfalse$
என்பதன் சுருக்கம்.

\end{document}
